% Page 080

\seite{80}

12.\ April.  Am 16.\ März fand eine Revision der Pfarr=\ob
schule Pfarrbezirk Mahrbach statt.\ob
Die Osterferien dauerten vom 31.\ März bis\ob
einschließlich 14.\ April.\ob
Aus der Schule entlassen wurden 4 Kinder,\ob
10 Knaben und 3 Mädchen. Neu aufgenommen\ob
wurden 4 Knaben und 2 Mädchen. Zu Beginn\ob
des neuen Schuljahres beträgt der Schülerstand\ob
42 Kinder, während im Vorjahr Differenz 21, es\ob
nur 34 Kinder die Schule besuchten.

22.\ Mai.    Mit dem heutigen Tag beginnen die Pfingstferien.\ob
Die Ferien bis 31.\ Mai einschließlich.

14.\ Juli.   Infolge der mangelhaften Witterungsverhältnisse wurde mit\ob
den Hundstagen nicht so früh begonnen wie im Vorjahre, sondern\ob
wurden die Ferienzeiten auf später gelegt worden. Es sind\ob
Hundsferien vom 4.– 14.\ Juli.\ob
21.\ August. Die diesjährigen Herbstferien dauern vom 19.\ September\ob
bis 21.\ Oktober.\ob
Ein großes Unglück ist vor längerer Zeit beim Dorf=\ob
brand von Bettingen über dieses Dorf herein gefallen.\ob
Die Häuser, die meist Stroh= und Strohdächer sind,\ob
werden in der Regel in einem feuermäßigen Umfange an=\ob
ge gewonnen, sodaß Menschen und Tiere Not ge=\ob
litten sind. Besonders sind es die Malberger Ordensleute,\ob
die zur Ordensschule ins Kloster fahren, die morgens\ob
und abends darin\unsicher{} abgefackelt, Viehlose u.\ Habern\unsicher\ob
in ihren Futterwägen machen. Oben\unsicher{} auf einmal\ob
am besten wurden am gewissen Sonntag zwei Pad=\ob
res von dem Seminar der Weißen Väter aus Trier\ob
entsandt. Der eine hielt eine Rede, dann in eine La=\ob
maschine zurück und fuhr Bett-Einsammeln\unsicher.

28.\ November.  Vom 16.– 20.\ November fiel der Unterricht infolge\ob
Krankheit des Lehrers Rings – (Halsleiden) – aus.\ob
Vom 22.– 27.\ November übernahm Lehrer Weiler\ob
von Olzheim den Unterricht.\ob
7.\ Dezember.  Am 23.\ November nahm Lehrer Rings den Unterricht\ob
wieder auf. Am 2.\ u. 3.\ Dezember mußte Unterricht da infolge\ob
der Krankheit des Lehrers\unsicher.
\pend
